\setcounter{deso}{3}
\begin{name}
	{\tenchude}
	{\tendethi}
	{\tentruong}
	{\thoigian}
\end{name}
\begin{bt}[1,5 điểm]~
\begin{enumerate}[1.]
    \item Tính giá trị biểu thức $A=\sqrt{(2-\sqrt{5})^2}-\sqrt{100}+\sqrt{45}$.
    \item Giải phương trình $5x-3=-2(x-3)$.
    \item Giải bất phương trình $2x-1 \ge 5-4x$.
\end{enumerate}
\end{bt}

\begin{bt}[2,0 điểm]~
\begin{enumerate}[1.]
    \item Cho hàm số $y=ax^2$ với $a \ne 0$ có đồ thị cắt đường thẳng $y=x$ tại hai điểm có hoành độ là $0$ và $1$. Xác định hệ số $a$.
    \item Cho biểu thức $M=\dfrac{x-2\sqrt{x}+1}{x-1}-\dfrac{1}{2\sqrt{x}+2}$ và $N=\dfrac{\sqrt{x}-1}{\sqrt{x}}$
    \begin{enumerate}
        \item Rút gọn biểu thức $M$.
        \item Tìm giá trị của $x$ để $M=2$.
    \end{enumerate}
\end{enumerate}
\end{bt}
\begin{bt}[1,5 điểm]%[9D4V2-5]
	Hai thành phố A và B cách nhau $200$ km. Một ô tô đi chuyển từ A đến B, rồi quay trở về A. Biết tốc độ lúc đi lớn hơn tốc độ lúc về là $10$ km/h. Do đó, thời gian về nhiều hơn thời gian đi là $1$ giờ. Tính tốc độ lúc đi của ô tô.
	\loigiai{
	Gọi tốc độ lúc đi của ô tô là $x$ (km/h; $x>10$).\\
	Vì tốc độ lúc đi lớn hơn tốc độ lúc về là $10$ km/h nên ta có tốc độ lúc về là $x-10$ (km/h).\\
	Thời gian ô tô đi từ A đến B là $\dfrac{200}{x}$ (giờ).\\
	Thời gian ô tô từ B trở về A là $\dfrac{200}{x-10}$ (giờ).\\
	Vì thời gian về nhiều hơn thời gian đi là $1$ giờ nên ta có phương trình:
	\begin{eqnarray*}
	&&\dfrac{200}{x-10} - \dfrac{200}{x} = 1\\
	&&\dfrac{200x}{x(x-10)} - \dfrac{200(x-10)}{x(x-10)} = \dfrac{x(x-10)}{x(x-10)}\\
	&&200x - 200(x-10) = x(x-10)\\
	&&200x - 200x + 2\,000 = x^2 - 10x\\
	&&x^2 - 10x - 2\,000 = 0 \quad(*)
\end{eqnarray*}
	Xét phương trình $(*)$ có $a=1$, $b=-10$, $c=-2\,000$:\\
	$\Delta = b^2 - 4ac = (-10)^2 - 4\cdot1\cdot(-2\,000) = 8\,100 > 0$.\\
	Do đó phương trình $(*)$ có hai nghiệm phân biệt:
	$$x_1 = \dfrac{-b+\sqrt{\Delta}}{2a} = \dfrac{-(-10)+\sqrt{8\,100}}{2\cdot1} = 50 \text{ (thỏa mãn).}$$
	$$x_2 = \dfrac{-b-\sqrt{\Delta}}{2a} = \dfrac{-(-10)-\sqrt{8\,100}}{2\cdot1} = -40 \text{ (loại).}$$
	Vậy tốc độ lúc đi của ô tô là $50$ km/h.
	}
\end{bt}

\begin{bt}[1,5 điểm]%[9D5H1-1]%[9D6H2-2]%
	\begin{enumerate}
\item Một cơ sở chăn nuôi gia cầm tiến hành nuôi thử nghiệm giống gà đẻ trứng mới. Khi gà đã cho trứng họ tiến hành khảo sát với $20$ quả trứng được cân nặng (gam) như sau:
	\begin{center}
\begin{tabular}{|c|c|c|c|c|c|c|c|c|c|}
	\hline
	$40$ & $42$ & $39$ & $38$ & $40$ & $42$ & $32$ & $40$ & $39$ & $38$ \\
	\hline
	$38$ & $40$ & $40$ & $40$ & $40$ & $39$ & $40$ & $39$ & $40$ & $42$ \\
	\hline
\end{tabular}
	\end{center}
Lập bảng tần số cho mẫu số liệu trên.
\item Trong trò chơi \lq\lq Chiếc nón kỳ diệu\rq\rq, khi người chơi quay ngẫu nhiên một lần, chiếc nón dừng lại tại một trong $19$ ô hình quạt, mỗi ô tương ứng là số điểm, trong đó có một số ô đặc biệt như hình bên và các ô có khả năng xảy ra như nhau. Hãy tính xác suất của biến cố $A\colon$\lq\lq Người chơi quay trúng ô $100$ điểm\rq\rq.
%%%%%%%%%%%%%%Hình vòng quay số
\begin{center}
	\begin{tikzpicture}[line join = round, line cap=round,>=stealth,font=\footnotesize,scale=0.6]
				\draw (0,0) circle (5cm);
	%	\draw[fill=black] (0,0)--(-110:6)--(-70:6)--cycle;
		\foreach \x in {90,109,128,147,166,185,204,223,242,261,280,299,318,337,356,375,394,413,432}{
			\coordinate (A\x) at (\x:5);
			\draw (0,0)--(\x:5)
			;
		}
			\foreach \m/\n/\i in {A90/A109/100,A128/A147/200,A147/A166/300,A185/A204/900,A204/A223/800,A242/A261/100,A261/A280/200,A299/A318/300,A318/A337/400,A356/A375/300,A375/A394/700,A394/A413/600,A432/A90/500}{
		\path(0,0)--($(\m) !0.5!(\n)$)node[pos=0.85,sloped]{$\i$};}
			\foreach \m/\n/\i in {A109/A128/Thêm lượt,A166/A185/Thưởng,A223/A242/Mất điểm,A280/A299/Nhân $2$,A337/A356/May mắn,A413/A432/Chia $2$}{
			\path(0,0)--($(\m) !0.5!(\n)$)node[pos=0.7,sloped]{\i};}
		\draw[fill=black] (0,0) circle (0.3cm);
		\draw[fill=white] (0,0) circle (0.2cm);
		%%%mũi tên
		\draw[fill=black,rotate=9] 
		(120:0.3) -- (95:1)--($(95:1)+(-0.3,0)$)--(90:2.5)--($(85:1)+(0.3,0)$)--(85:1)--(60:0.3)
		;
	\end{tikzpicture}
\end{center}
%%%%%%%%%%%%
	\end{enumerate}
	\loigiai{
		\begin{enumerate}
			\item Bảng tần số cho mẫu số liệu trên
				\begin{center}
				\begin{tabular}{|c|c|c|c|c|c|}
					\hline
					Cân nặng (gam) &$32$&$38$&$39$&$40$&$42$ \\
					\hline
					Tần số &$1$&$3$&$4$&$8$&$4$ \\
					\hline
				\end{tabular}
				\end{center}
				\item Có $19$ kết quả có thể khi quay ngẫu nhiên chiếc nón.\\
				Có $2$ ô $100$ điểm nên có $2$ kết quả thuận lợi cho biến cố $A$.\\
				Vậy xác suất của biến cố $A\colon$\lq\lq Người chơi quay trúng ô $100$ điểm\rq\rq~là $\dfrac{2}{19}$.
		\end{enumerate}
	}
\end{bt}
%%%==============HetBai_BT6==============%%%

%%%==============Bai_BT7==============%%%
\begin{bt}[1,0 điểm]%[Đề tham khảo tuyển sinh vào 10 - Quận 3 - Đề số 3]%[TS10Y6]
\immini{
Nước giải khát thường đựng trong lon nhôm và cỡ lon phổ biến chứa được khoảng $330$ ml chất lỏng, được thiết kế hình trụ với chiều cao khoảng $10{,}2$ cm (phần chứa chất lỏng), đường kính đáy khoảng $6{,}42$ cm. Nhưng hiện nay các nhà sản xuất có xu hướng tạo ra những lon nhôm với kiểu dáng cao thon hơn. Tuy chi phí sản xuất những chiếc lon cao này tốn kém hơn, nhưng nó lại dễ đánh lừa thị giác và được người tiêu dùng ưa chuộng hơn. }
{
        \begin{tikzpicture}[>=stealth, line join=round, line cap=round, font=\footnotesize, scale=1]
            \def\htru(#1,#2)(#3){
                \def\x{#1} % Bán kính trụ lớn
                \pgfmathsetmacro{\y}{\x/4} % Bán kính trục bé
                \def\h{#2} % Chiều cao
                \draw ($(#3)+(0:\x)$) arc (0:-180:{\x} and {\y})--++(90:#2) arc (180:0:{\x} and {\y}) arc (0:-180:{\x} and {\y}) ($(#3)+(0:\x)$)--++(90:\h);
                \draw[dashed]
                ($(#3)+(0:\x)$) arc (0:180:{\x} and {\y});
            }%%%
            \path
            (0,0) coordinate (O)
            (4,0) coordinate (A);
            \htru(2,2)(O)
            \htru(1,4)(A)
%            \path
%            (O)--(0:2)node[above,midway]{$2R$}
%            ($(O)+(180:2)$)--++(90:2)node[midway,left]{$a$}
%            (A)--++(0:1)node[midway,above]{$R$}
%            ($(A)+(180:1)$)--++(90:4)node[midway,left]{$2a$}
%            ;
        \end{tikzpicture}
}
\begin{listEX}
\item Một lon nước ngọt cao $13,41$ cm (phần chứa chất lỏng), đường kính đáy là $5,6$ cm. Hỏi lon nước ngọt cao này có thể chứa được hết lượng nước ngọt của một lon có cỡ phổ biến không? Vì sao?
\begin{center}
Biết thể tích hình trụ: $V=\pi r^2 h$ với $\pi \approx 3,14$
\end{center}
\item Biết chi phí sản xuất một chiếc lon tỉ lệ thuận với diện tích toàn phần của lon. Hỏi chi phí sản xuất chiếc lon cao tăng bao nhiêu phần trăm so với chi phí sản xuất chiếc lon cỡ phổ biến? (làm tròn $1$ chữ số thập phân).\dapso{$5{,}4\%$}
\begin{center}
Biết diện tích xung quanh, diện tích toàn phần hình trụ được tính theo công thức:
$S_{xq}=2\pi r h$ và $S_{tp}=S_{xq}+2S_{\textit{đáy}}$
\end{center}
\end{listEX}
\loigiai{
\begin{listEX}
\item
Thể tích chính xác của lon nhôm phổ biến là
$$V=\pi R^2 h=3{,}14\left(\dfrac{6{,}42}{2}\right)^2 10{,}2\approx 330{,}0198 \text{ cm}^3.$$
Thể tích của lon nhôm cao là
$$V=\pi R^2 h=3{,}14\left(\dfrac{5{,}6}{2}\right)^2 13{,}41\approx 330{,}122 \text{ cm}^3.$$
Vậy lon nước cao có thể chứa được hết lượng nướccủa lon phổ biến.
\item Diện tích xung quanh của lon phổ biến là
$$S_{xq}=2\pi r h=2\cdot 3{,}14 \cdot\dfrac{6{,}42}{2}\cdot 10{,}2 \approx 205{,}62 \text{ cm}^2.$$
Diện tích đáy của lon phổ biến là $$S_{\textit{đáy}}=\pi r^2 =3{,14}\left(\dfrac{6{,}42}{2}\right)^2\approx 32{,}35 \text{ cm}^2.$$
Diện tích toàn phần của lon phổ biến là $$S_{tp}=S_{xq}+2S_{\textit{đáy}}=205{,}62 +2\cdot 32{,}35 =270{,}32\text{ cm}^2.$$
Diện tích toàn phần của lon nhôm cao là $$S'_{tp}=S_{xq}+2S_{\textit{đáy}}=2\cdot 3{,}14 \cdot\dfrac{5,6}{2}\cdot 13,41 +2\cdot 3{,14}\left(\dfrac{5,6}{2}\right)^2 \approx 285{,}03\text{ cm}^2.$$
Vậy diện tích toàn phần của lon nhôm cao bằng $\dfrac{285{,}03}{270{,}32}\cdot 100\%\approx 105{,}4\%$ so với lon nhôm phổ biến.\\
Do đó chi phí sản xuất chiếc lon cao tăng $5{,}4\%$ so với chi phí sản xuất chiếc lon cỡ phổ biến.
\end{listEX}
}
\end{bt}

\begin{bt}[2,5 điểm]%[9H3V4-1]
	Cho đường tròn $(O;R)$ và một điểm $S$ nằm ngoài đường tròn. Từ điểm $S$ kẻ hai tiếp tuyến $SA$, $SB$ với đường tròn $(O;R)$ ($A$, $B$ là các tiếp điểm).
	\begin{enumerate}
		\item Chứng minh tứ giác $O A S B$ là tứ giác nội tiếp.
		\item Kẻ đường kính $BD$ của đường tròn $(O;R)$. Đường thẳng $SD$ cắt đường tròn $(O;R)$ tại $C$ ($C$ khác $D$). Gọi $I$ là giao điểm của $S O$ và $A B$. Tia $CI$ cắt đường tròn $(O;R)$ tại điểm thứ hai là $M$. Chứng minh $\triangle SCI$ đồng dạng với $\triangle SOD$ và $SO$ song song với $BM$.
	\end{enumerate}
	\loigiai{
		\begin{center}
			\begin{tikzpicture}[scale=0.8, font=\footnotesize, line join=round, line cap=round, >=stealth]
				\def\r{3}
				\path (0,0) coordinate (O)
				(125:\r) coordinate (A)
				(-125:\r) coordinate (B)
				($(A)!1mm!90:(O)$) coordinate (x)
				($(B)!1mm!90:(O)$) coordinate (y)
				(intersection of A--x and B--y) coordinate (S)
				(55:\r) coordinate (D)
				(-55:\r) coordinate (M);
				
				\draw[name path=c1] (O) circle (\r);
				\path[name path=dt] (S)--(D);
				\path[name intersections={of=dt and c1}] 
				(intersection-2) coordinate (C)
				($(S)!(A)!(O)$) coordinate (I)
				%				(intersection of A--D and C--B) coordinate (F)			
				;
				\draw 
				(S)--(D) (S)--(A)--(C)
				(S)--(B)--(A) (B)--(M)--(C) (S)--(O) (A)--(O)--(B) (A)--(O)--(D)--cycle (C)--(O)
				;
				\foreach \d/\g in {S/-180,O/-70,B/-90,C/120,A/90,D/60,M/-90,I/-120}
				\path[draw,fill=black] (\d) circle(1pt) + (\g:9pt) node {$\d$};
				\foreach \a/\b/\c in {S/A/O,S/B/O,A/I/O}{\draw pic[draw,angle radius=2mm] {right angle = \a--\b--\c};}
			\end{tikzpicture}
		\end{center}
		\begin{enumerate}
			\item Ta có $\triangle S A O$ vuông tại $A$ (do $S A$ là tiếp tuyến của $(O)$).\\
			Do đó $S$, $A$, $O$ cùng thuộc đường tròn đường kính $SO$. (1)\\ Lại có $\triangle SBO$ vuông tại $B$ (do $S B$ là tiếp tuyến của $(O)$).\\
			Do đó $S$, $B$, $O$ cùng thuộc đường tròn đường kính $SO$. (2)\\
			Từ (1) và (2) suy ra $S$, $A$, $B$, $O$ cùng thuộc đường tròn đường kính $SO$.\\
			Vậy $OASB$ là tứ giác nội tiếp.
			\item Ta có $\widehat{SAC}+\widehat{OAC}=\widehat{SAO}=90^{\circ}$ nên $\widehat{SAC}=90^{\circ}-\widehat{OAC}$. (3)\\
			$\triangle O A C$ cân tại $O$ (do $OA=OC$) nên
			$\widehat{OAC}=\dfrac{180^{\circ}-\widehat{AOC}}{2}=90^{\circ}-\dfrac{\widehat{AOC}}{2}=90^{\circ}-\widehat{ADC}$. (4)\\
			Từ (3) và (4) ta có $\widehat{S A C}=\widehat{S D A}$.\\
			Xét $\triangle S A C$ và $\triangle S D A$, ta có\\
			$\widehat{DSA}$ chung và $\widehat{SAC}=\widehat{SDA}$ (chứng minh trên).\\
			Do đó $\triangle SAC \backsim \triangle SDA$ (góc - góc).\\
			Suy ra $\dfrac{S A}{S C}=\dfrac{S D}{S A}$ hay $S A^2=S C \cdot S D$\\
			Do đó $SC \cdot SD=SI \cdot SO$ hay $\dfrac{SC}{SO}=\dfrac{SI}{SD}$.\\
			Xét $\triangle SCI$ và $\triangle SOD$, ta có\\
			$\widehat{D S O}$ chung và $\dfrac{S C}{S O}=\dfrac{S I}{S D}$.\\
			Do đó $\triangle S C I \backsim \triangle S O D$ (cạnh - góc - cạnh).\\
			Suy ra $\widehat{S I C}=\widehat{S D O}$.\\
			Mà $\widehat{S D O}=\widehat{C M B}$ (cùng chắn cung $B C$).\\
			Nên $\widehat{S I C}=\widehat{C M B}$.\\
			Hơn nữa hai góc này ở vị trí đồng vị nên $S I \parallel B M$.\\
			Vậy $S O \parallel B M$.
		\end{enumerate}
	}
\end{bt}